\section{Background and Motivation}
\label{sec_background}


\begin{table}[h]
    \centering
    \resizebox{\columnwidth}{!}{
    \begin{tabular}{|l|c|cccc|}
        \hline
        \multirow{2}{*}{\textbf{Model}} & \textbf{Context Window} & \multicolumn{4}{c|}{\textbf{Acc. at Input Len. (Tokens)~\cite{levy2024same}}} \\
        \cline{3-6}
         & \textbf{(Tokens)} & 250 & 1K & 2K & 3K \\
        \hline
        GPT-3.5~\cite{openai2022gpt35} & 4K & 0.72 & 0.64 & 0.60 & 0.59 \\
        GPT-4o~\cite{openai2023gpt4} & 128K & 0.95 & 0.92 & 0.86 & 0.80 \\
        Gemini Pro~\cite{google2024gemini} & 1M & 0.94 & 0.80 & 0.66 & 0.62 \\
        Mixtral 8x7B~\cite{mistral2023mixtral} & 32K & 0.90 & 0.72 & 0.68 & 0.60 \\
        \hline
    \end{tabular}
    }
\caption{Comparison of context window sizes and accuracy degradation with increasing input/prompt length.}
    \label{tab:context_accuracy}
\end{table}
% \aaron{The paragraphs below (leading up to Section 2.1) repeat much of the first half of the introduction. I propose we eliminate the details about model checkers and consistency checkers, and include details as appropriate in Section 2.1, where we compare LLMs with these prior tools. Furthermore, I propose we trim and combine the first and fourth paragraphs below to provide a concise lead-in to Section 2 that builds on what was already said in the introduction.}

% Routing device misconfigurations are a common and persistent issue, often leading to
% security vulnerabilities, performance degradation, or even network outages
% \cite{zheng2012atpg, feamster2005detecting}. Addressing these misconfigurations
% manually is akin to having a domain expert, such as a network engineer,
% painstakingly analyze every configuration line. While an expert leverages years
% of experience and an in-depth understanding of standard practices, the complexity
% and interdependency of modern configurations \cite{le2007rr, benson2009complexitymetrics}
% make this process not only labor-intensive but also prone to error—especially at scale.

% To reduce this reliance on manual effort, automated tools such as model checkers
% and consistency checkers have emerged. Model checkers
% \cite{fogel2015general, beckett2017general, abhashkumar2020tiramisu, prabhu2020plankton,
% 	zhang2022sre, steffen2020netdice, ye2020hoyan, ritchey2000using, al2011configchecker,
% 	jeffrey2009model} identify misconfigurations by creating a logical model of the
% control and data planes, then analyzing or simulating whether engineer-specified
% forwarding policies are satisfied. For example, Batfish simulates the control and data
% planes under specific network conditions (\eg, link failures) and checks whether
% forwarding paths deviate from specified policies \cite{fogel2015general}.
% \aaron{Categorize Batfish as consistency checker as well; Batfish does a form of comparing with best-practices}
% Similarly, Minesweeper uses a satisfiability solver to verify whether reachability
% and load-balancing policies hold under various scenarios (\eg, all single link
% failures) \cite{beckett2017general}. Although model checkers are popular, their
% efficacy depends heavily on the correctness and completeness of the predefined
% policies, which in turn require significant domain expertise to design. Moreover,
% they often struggle with the nuances of particular protocols and vendor-specific
% features \cite{ye2020hoyan, birkner2021metha}.

% In contrast, consistency checkers such as \textit{Diffy} \cite{kakarla2024diffy}
% and \textit{Minerals} \cite{le2006minerals} use statistical methods to learn
% typical configuration patterns, identifying deviations as potential errors rather
% than relying on predefined policies. For instance, Diffy learns configuration
% templates to spot anomalies in new configurations, while Minerals applies
% association rule mining to detect misconfigurations by deriving local policies
% from router files. Unfortunately, these approaches often oversimplify configurations
% by treating any deviation from standard patterns as misconfigurations—leading
% to false positives whenever valid, network-specific or device-specific choices
% differ from general norms.

As mentioned previously, many automated tools have been developed to combat misconfigurations on routers and switches.
%To efficiently combat misconfigurations on routers and switches, many automated tools 
%such as model checkers and consistency checkers 
%have been developed. 
Model checkers, such as Batfish~\cite{fogel2015general, batfish}, Minesweeper~\cite{beckett2017general}, and Tiramisu~\cite{abhashkumar2020tiramisu}, create logical models of network behavior to verify compliance with predefined policies. However, their reliance on manually specified policies limits scalability and adaptability~\cite{ye2020hoyan, birkner2021metha}. Consistency checkers, such as Batfish~\cite{batfish},\footnote{Batfish~\cite{batfish} can function as a model checker or a consistency checker, depending on which capabilities are invoked.} Diffy~\cite{kakarla2024diffy}, and Minerals~\cite{le2006minerals}, identify deviations from learned patterns or best practices, offering a more automated approach. However, these tools often oversimplify configurations, flagging valid variations as errors and leading to false positives. More recently, Large Language Models (LLMs)~\cite{bogdanov2024leveraging, chen2024automatic, wang2024identifying} have been explored as a promising alternative. Pretrained on vast network-related datasets, LLMs combine semantic understanding with flexibility, enabling them to reason about both local and global dependencies in configurations.
Unlike consistency checkers, LLMs adapt to network contexts and detect nuanced misconfigurations, and unlike model checkers, they automate detection without requiring manually defined rules, as demonstrated by tools like Ciri~\cite{lian2023configuration}.
A notable example is
Ciri~\cite{lian2023configuration}, an LLM-based configuration verifier that
illustrates the potential for scalable and accurate automated detection.

    
% Given the limitations of existing tools, an ideal solution for misconfiguration
% detection would employ an automated, intelligent \emph{agent} with expert-level
% knowledge of routing device configurations. Such an agent could examine configuration
% lines in detail and apply domain expertise to detect potential issues. Recent
% advances in Large Language Models
% (LLMs)~\cite{bogdanov2024leveraging,chen2024automatic,wang2024identifying,liu2024large,
% 	wang2024netconfeval} provide a promising avenue toward this goal. Pre-trained on
% large volumes of network-related data, these models embed a deep understanding of
% best practices and standards, enabling them to analyze intricate configurations
% effectively. At the same time, LLMs' successful applications across diverse
% domains~\cite{carion2020end,sheng2019nrtr,neil2020transformers,parmar2018image,
% 	chen2021developing,gulati2020conformer} make their use in network misconfiguration
% detection a logical next step. A notable example is
% Ciri~\cite{lian2023configuration}, an LLM-based configuration verifier that
% illustrates the potential for scalable and accurate automated detection.


\subsection{How LLMs Encode and Apply Domain Knowledge}

\mysubpara{Encoding Network Expertise:} LLMs, particularly transformer-based generative pre-trained models~\cite{vaswani2017attention,hill2024transformers,lin2022survey}, encode domain knowledge during pretraining by processing extensive datasets of network documentation, configuration files, and protocol specifications~\cite{qiu2020pre,achiam2023gpt,touvron2023llama,brown2020language}. Through this process, the models learn statistical patterns, structural hierarchies, and semantic relationships crucial to network operations. For example, they identify how ACL rules typically interact with routing policies or how specific firewall configurations align with best practices. This knowledge is embedded within the model’s parameters, which are refined using gradient descent to minimize prediction errors, effectively transforming the model into a probabilistic repository of domain-specific expertise.
Technically, transformers achieve this by employing multi-head self-attention mechanisms, which dynamically assign importance to critical tokens such as protocol keywords or IP ranges, and positional encodings, which preserve the sequence of commands. These mechanisms allow the model to capture both surface-level syntax (\eg, ACL formats) and deeper semantic relationships (\eg, precedence of routing policies over firewall rules). The resulting parameterized knowledge acts like a domain expert’s internalized understanding, enabling the model to generalize its expertise to new configurations it encounters.

\mysubpara{Advantages Over Non-LLM-Based Tools:} The unique architecture of transformer-based LLMs allows them to fundamentally surpass consistency checkers by integrating pre-trained domain knowledge with context-specific information from the configuration being analyzed, creating a much richer decision-making framework. Consistency checkers rely solely on statistical patterns derived from historical or similar data, using majority-voting logic to flag deviations as anomalies. This approach assumes that frequently observed patterns are always correct and treats all deviations as errors, regardless of whether they represent valid, context-specific configurations. As a result, consistency checkers fail to account for the underlying intent of the configuration or the broader network context, often leading to false positives and missed opportunities to understand nuanced misconfigurations.
In contrast, LLMs integrate their pre-learned understanding of best practices with the specific configuration context provided in the prompt. By treating configurations as interconnected systems, they can assess whether deviations align with learned standards or represent genuine misconfigurations. For instance, while a consistency checker might flag a rare ACL rule as an anomaly, an LLM can reason that it is valid based on its relationships to routing policies or protocol directives elsewhere in the configuration.

Compared to model checkers, LLMs are fully automated, requiring no human-defined rules and dynamically adapting to new contexts, making them better suited for handling configurations involving vendor-specific nuances or unforeseen patterns. In contrast, model checkers simulate network behavior and verify compliance with predefined policies but rely heavily on the manual specification of rules, algorithms, and policies, which limits their scalability due to dependence on domain expertise. However, LLMs have certain limitations when compared to model checkers, particularly in their inability to rigorously simulate and compute precise route behaviors under specific scenarios, such as link failures or traffic balancing. While model checkers can mathematically verify properties like load balancing, LLMs rely on probabilistic reasoning and can only infer whether a configuration aligns with best practices or observed patterns. Thus, while LLMs excel in automation and contextual reasoning, they lack the exact computational guarantees offered by model checkers. The focus of this paper, however, is not to address the mathematical computational limitations of LLMs but to optimize existing LLMs for effective and scalable misconfiguration detection.

\mysubpara{How LLMs Integrate Prompts and Pretrained Knowledge for Context-Aware Reasoning:} When a user provides configuration snippets and instructions in a prompt (\eg, “Evaluate this ACL for conflicts”), the LLM combines its pre-trained knowledge with the contextual information in the prompt. The self-attention mechanism enables the model to cross-reference tokens in the prompt with its stored knowledge, such as identifying a conflict between a firewall rule and routing policy based on probabilistic relationships learned during pretraining. By estimating the likelihood of each token or sequence (\(p(\text{token}_i \mid \text{token}_1, \ldots, \text{token}_{i-1})\)), the model highlights deviations from high-probability patterns, flagging potential misconfigurations such as overlapping IP ranges or misplaced \texttt{deny} directives. This integration mirrors how a human domain expert draws on years of experience to evaluate new configurations, applying abstract principles to specific cases.


\subsection{Building Context in Prompts for Accurate Misconfiguration Detection}

Through this architecture and the pre-learned general network configuration
knowledge, LLM-based Q\&A models detect subtle errors that can emerge from
interactions among different elements—issues frequently missed by conventional
tools. As networks grow more complex and interdependent, the model’s ability to
holistically evaluate configurations and understand their intended functions
makes it a valuable asset for enhancing misconfiguration detection accuracy.
However, to realize this capability, query prompts must include all relevant
context from the configuration file when analyzing particular configuration lines
or excerpts, including associated segments that might reveal hidden
misconfigurations. Incorporating this broader view of the file refines the
model’s response by leveraging network-specific context unique to the active
configuration, rather than relying solely on pre-learned knowledge. Without
this detailed context, the LLM’s capacity to identify potential issues is
significantly diminished~\cite{liskavets2024prompt,tian2024examining,khurana2024and,
	shvartzshnaider2024llm}.

\mypara{Limitations of Full-File Prompting:}
A straightforward solution to providing extensive context is to feed the entire
configuration file to the LLM—either as a single prompt or progressively using chain of thoughs (CoT)—and let
the model detect potential misconfigurations. Yet, this approach is highly inefficient
due to token-length constraints~\cite{xue2024repeat,yu2024breaking,gu2023mamba}.
Moreover, it can trigger \emph{context overload}~\cite{lican,li2024long,qian2024long},
where extraneous information distracts the model from the crucial lines in the file,
leading to missed errors or false positives. Analogously, a human expert would likely
be overwhelmed by an avalanche of irrelevant details, making it harder to pinpoint
actual issues. Similarly, context overload diminishes the model’s ability to focus
and raises the risk of missing important misconfigurations or erroneously flagging
innocuous lines.
Table~\ref{tab:context_accuracy} underscores a key insight: while LLMs with larger context windows (\eg, Gemini Pro’s 1M tokens) can process entire configuration files, accuracy degrades as input length increases. Even models with smaller windows, such as GPT-3.5 (4K tokens) and Mixtral 8x7B (32K tokens), show similar trends—accuracy drops from 250 to 3,000 tokens, with CoT offering little improvement. This issue is critical in real-world router configurations, where even moderate-sized Enterprise Edge Router files (15K–50K tokens) exceed the reliable reasoning limits of most models. Carrier-grade routers (100K–500K tokens) far surpass this threshold, making full-file prompting ineffective regardless of context size. Ultimately, the challenge is not just context length but the model’s ability to sustain accurate reasoning over large inputs, where full-file prompting—whether in one shot or iteratively—yields diminishing returns.

\mypara{Challenges of Partition-Based Prompting:}
A common solution to context overload has been partition-based prompting,
where the large and complex configuration file is broken down into smaller
chunks or sections, typically based on the order in which the configuration
appears in the file, such as Ciri~\cite{lian2023configuration}. 
 Although this method eases token constraints, it also fragments
vital context spread across multiple sections. As network configurations often include
parameters that interact or depend on other areas of the file, merely analyzing each
partition in isolation may overlook key interdependencies. In contrast, a proficient
human engineer would cross-reference earlier and later sections continuously; however,
a purely partitioned prompting approach discards this global perspective. While prompt
chaining~\cite{wang2024identifying,bogdanov2024leveraging} can maintain some continuity
by linking prompts together, it does not inherently address “context chaining”, and
thus crucial interactions between file sections can remain undetected.

Consider a file where firewall rules in one segment conflict with routing policies
in another. Partition-based prompting would handle these sections separately, missing
the cross-sectional conflict entirely. This fragmentation of context is particularly
troublesome for dependency-related misconfigurations, which demand an integrated view
of how distinct segments interrelate. If a given chunk of text solely contains lines
on unrelated parameters, context about relevant policy sections is lost, rendering
the model unable to draw necessary conclusions.

In this paper, we introduce a tool that tackles these issues via a \emph{context-aware,
	iterative prompting mechanism}. This approach strategically manages and incorporates
only the relevant parts of the configuration for each line under review, ensuring that
LLMs obtain the required context without succumbing to overload. In the next section,
we discuss the specific challenges encountered during the development of this solution
and outline our methodology in detail.

